% Verbatim copy of /Users/alingalatan/workspaces/grothe/handoff/manuscript/legacy/lower-bound/sharp-circular.tex (file dated 2026-10-25). Only this line was added.
\subsection{A nearly optimal circular-mean slope and a larger shared-A region}
\label{sec:sharp-circular-audit}
\paragraph{Status and scope.}
\st{P,C}: the analytic reduction and directed constants below are new.
The conclusions about $S$ additionally use the already supplied benchmark
$S-B_0>0.0020476$ from \S\ref{sec:sar-second}; the package also replays the unchanged inherited Jensen benchmark
checker separately from the new Python checks. No full K or endpoint is claimed.

\begin{theorem}[Global circular-mean majorant]\label{thm:sharp-circular-audit}
For $m(s)=(2\pi)^{-1}\int_0^{2\pi}|1+\sqrt{s}e^{i\theta}|\,d\theta$,
\[
 m(s)\le1+\frac{59767}{200000}s\qquad(s\ge0).
\]
If $C_* =\sup_{s>0}(m(s)-1)/s$, then
$0.298834<C_*\le0.298835$. Thus the coefficient is within $10^{-6}$
of the best possible global affine coefficient with intercept one.
\end{theorem}
\begin{proof}
Put $a_k=\binom{1/2}{k}^2$. For $0\le s\le1$,
$m(s)=\sum_{k\ge0}a_ks^k$ and $\sum a_k=4/\pi$.
The coefficients are nonnegative, so
$m(s)\le1+(4/\pi-1)s$, which is stronger than claimed.
For $s=r^2\ge1$, inversion gives $m(r^2)=r m(r^{-2})$.
With $N=12$ and
$d_N=31831/25000-\sum_{k=0}^N a_k>0$, the exact inequality
$4/\pi<31831/25000$ gives
\[
 m(r^2)\le\sum_{k=0}^N a_k r^{1-2k}+d_Nr^{-2N-1}.
\]
Writing $C=59767/200000$, it suffices to show positivity of the rational
polynomial
\[
 P(r)=Cr^{2N+3}-r^{2N+2}+r^{2N+1}
       -\sum_{k=1}^N a_kr^{2N+2-2k}-d_N.
\]
The replay checks positive Bernstein coefficients on eight rational
subintervals covering $[1,3]$, and positive coefficients of $P(3+x)$,
covering $[3,\infty)$. These finite comparisons prove the global upper
bound. For the lower bound on $C_*$, use $r=581/400$ and the positive
partial sum $r\sum_{k=0}^{24}a_kr^{-2k}$; exact arithmetic gives
$(r\sum_{k=0}^{24}a_kr^{-2k}-1)/r^2>149417/500000$.
\end{proof}

\begin{theorem}[Orthogonality-only shared-A bound]\label{thm:shared-a-audit}
Let $q,r$ be any jointly distributed complex square-integrable variables,
independent of the visible standard proper complex Gaussian $Z$, with
$\E[\bar q r]=0$. No independence between $q$ and $r$, no Gaussian law
for them, and no chaos-degree assumption are needed. At the peak profile,
put $u=\E|q|^2$, $w=\E|r|^2$, $T=|Z|^2$, and
$B_0=2\E|L|=(387/500)\sqrt\pi$. Then
\begin{equation}\label{eq:shared-a-audit}
 \boxed{\Lambda(q+(T-1)r)
       \le B_0+0.38968u-0.12432w.}
\end{equation}
All terminating decimals here are exact rational bounds.
\end{theorem}
\begin{proof}
Averaging the independent phase of $Z$ gives exactly
$2\E[\ell(T)m(|q+(T-1)r|^2/\ell(T)^2)]$, where
$\ell(t)=\sqrt t(a+bt)$, $a=699/1250$, $b=179/1250$.
The variance penalty is $p(u+w)$, with $p=2500/1891$.
Orthogonality and the preceding majorant therefore give
\[
 \Lambda\le B_0+(2CK-p)u+(2CJ-p)w,
 \qquad K=\E\ell(T)^{-1},\quad J=\E (T-1)^2/\ell(T).
\]
Both integrals have closed forms. With $z=a/b$,
\[
 K=\frac{\pi e^z\operatorname{erfc}(\sqrt z)}{\sqrt{ab}},\quad
 I_1=\frac{\sqrt\pi-aK}{b},\quad
 I_2=\frac{\sqrt\pi/2-aI_1}{b},\quad J=I_2-2I_1+K.
\]
The recurrence follows by dividing $t$ by $a+bt$ in the integrals;
the first formula follows from the Gaussian resolvent integral.
Directed arithmetic proves $K<2.864$ and $J<2.004$, hence
$2CK-p<0.38968$ and $2CJ-p<-0.12432$.
The apparent singularity at $T=0$ is integrable; all expectations are
finite under the stated second-moment assumptions.
\end{proof}

\paragraph{Consequences for cubic shared drivers.}
For remote $q\in H_{2,1}$ and any shared linear $r\in H_{1,0}$,
orthogonality is automatic. With the inherited benchmark gap, the new
bound gives
\[
 \boxed{u\le0.00525+0.319w\ \Longrightarrow\ \Lambda<S,}
 \qquad
 \boxed{u\le0.00551+0.319w\ \Longrightarrow\ \Lambda<S+10^{-4}.}
\]
These improve the earlier $0.00541+0.292w$ weak region. Other visible
sectors are not present and cannot be appended without a simultaneous comparison.

\begin{theorem}[Larger shared-A strips, \st{P,C}]\label{thm:shared-a-strip-audit}
Under the same orthogonality and visible-independence assumptions, for
$0\le u\le0.0463$ the following implications hold:
\[
\begin{array}{c|c}
 w\ge0.0228&\Lambda<S-10^{-4}\\
 w\ge0.0225&\Lambda<S\\
 w\ge0.0223&\Lambda<S+10^{-4}.
\end{array}
\]
They use the same inherited lower bound $S-B_0>0.0020476$.
\end{theorem}
\begin{proof}
Let $C=59767/200000$ and
$g_C(s)=\min(1+Cs,\sqrt{1+s})$. This is an increasing concave majorant
of $m$, with $0\le g_C'\le C$: the nonzero crossing is
$s_0=(1-2C)/C^2$ and the derivative drops from $C$ to
$C/(2(1-C))$. Jensen therefore bounds $\Lambda$ by
\[
 V_C(u,w)=2\E\ell(T)g_C\!\left(\frac{u+(T-1)^2w}{\ell(T)^2}\right)-p(u+w).
\]
Its increments in $w$ are at most $(2CJ-p)\,\Delta w<0$.
For each boundary value $w_0$, use the square-root branch on
$0\le T\le d$ and the affine branch elsewhere, where
$r=11/50$ and $d=r^2$. This fixed choice is an upper bound $W(u,w_0)$
even when it does not choose the smaller branch. Writing
$h(x)=a+bx^2$ and $y(x)=u+w_0(x^2-1)^2$ gives
\begin{align*}
 W(u,w_0)-B_0={}&(2CK-p)u+(2CJ-p)w_0\\
 &+\int_0^r e^{-x^2}\left[
 4x\left(\sqrt{x^2h(x)^2+y(x)}-xh(x)\right)
 -\frac{4Cy(x)}{h(x)}\right]dx.
\end{align*}
There is no truncated tail: the affine part is integrated by the exact
$K,J$ formulas. On $0\le u\le u_0=0.0463$,
\[
 \partial_uW\ge
 \frac{de^{-d}}{\sqrt{d(a+bd)^2+u_0+w_0}}
 +2C\left(\frac{\sqrt\pi}{a+b/2}-\frac{2r}{a}\right)-p>0.
\]
Here $K\ge\sqrt\pi/(a+b/2)$ is Jensen under the
$\operatorname{Gamma}(1/2)$ law. Thus it suffices to evaluate
$W(u_0,w_0)$. Outward integer intervals on 8192 equal cells of $[0,r]$
prove, respectively,
\[
 W(u_0,0.0228)-B_0<0.001904,\quad
 W(u_0,0.0225)-B_0<0.002026,\quad
 W(u_0,0.0223)-B_0<0.002107.
\]
The verifier bounds $e^{-x^2}$ between its degree-nine and degree-eight
alternating Taylor polynomials. All square roots are enclosed by integer
square roots, and it also checks the three positive derivative bounds.
Monotonicity of $V_C$ in $w$, followed by $V_C\le W$, proves each entire
unbounded strip, not only its boundary.
\end{proof}

Combining the weak strip with the affine cone, the remaining part of
$u\le0.0463$, $w\le0.0328$ is contained in
\[
 w<0.0223,\qquad u>0.00551+0.319w.
\]
This conclusion still concerns the shared-A family only.

The nearly sharp slope also isolates the limitation of this particular
method: further substantial improvement requires an envelope retaining
more than the second moment, or the original-response constraints. Merely
optimizing the global affine slope has less than $10^{-6}$ left to gain.
